\documentclass{situationist-preprint}

\title{The Loop Variant: A Separate Proof}
\msid{ETH-03-c}
\msfamily{ETH-03}
\msdiscipline{Normative and Applied Ethics}
\mstype{Alternative proof}
\msstatus{Revised 7 October 2026}
\mskeywords{trolley problem; Loop case; doctrine of double effect; redirected threat}

\begin{document}
\maketitle

\begin{abstract}
In Thomson's Loop case, the side track loops back to the main track, so that the diverted trolley would kill the five from the other side unless stopped by the body of the one person on the loop. We give a proof, independent of the uniqueness theorem of ETH-03-a, that diverting the trolley in Loop is permissible. The proof uses only the doctrine of the redirected threat and a principle about the irrelevance of track geometry, and it makes no appeal to the doctrine of double effect.
\end{abstract}

\paragraph{Note on the revision.} The version of 6 October 2026 argued in Section~3 that the one's death in Loop is not intended, because the agent would divert the trolley even if the one were not on the loop. This is an appeal to the doctrine of double effect, which the family is meant to avoid, and it is moreover doubtful as a description of the case, since without the one the diversion would not save the five. Section~3 has been rewritten so that the argument relies only on the doctrine of the redirected threat.

\section{The case}

\citet{Thomson1985} introduced Loop to test the doctrine of double effect. In Switch, the one's death is a side effect of diverting the trolley; in Loop, the diversion saves the five only because the trolley hits the one, so that his being hit is a means to their survival. The doctrine of double effect therefore forbids diverting in Loop, while many judge it permissible.

\section{Two principles}

\begin{principle}[Redirection]\label{p:drt}
It is permissible to redirect an existing threat so that it kills fewer people than it would otherwise kill.
\end{principle}

\begin{principle}[Geometry]\label{p:geom}
Whether an act of redirection is permissible does not depend on the shape of the track beyond the point at which the redirected threat reaches its victims.
\end{principle}

Principle~\ref{p:drt} is the doctrine of the redirected threat (ETH-03-a). Principle~\ref{p:geom} is a special case of the invariance condition of ETH-03-a, stated for the variations that distinguish Loop from Switch.

\section{The proof}

\begin{theorem}\label{t:loop}
Diverting the trolley in Loop is permissible.
\end{theorem}

\begin{proof}
In Loop, the trolley is an existing threat to the five. Diverting it onto the loop redirects that threat, and the redirected threat kills one person, the one on the loop, rather than five. By Principle~\ref{p:drt}, the diversion is permissible, unless something about Loop distinguishes it from Switch. Loop differs from Switch only in the shape of the track beyond the point at which the trolley reaches the one: in Switch the side track continues away from the five, in Loop it rejoins the main track. By Principle~\ref{p:geom}, this difference does not affect permissibility. Hence the diversion in Loop is permissible.
\end{proof}

The proof does not ask whether the one's death is intended, foreseen, or a means. It asks only whether the act redirects an existing threat so that fewer are killed, and whether Loop differs from Switch in a way that matters. That the one's body stops the trolley is a fact about what happens after the threat has been redirected onto him; it does not make the act one of introducing a new threat, since the one is in the path of the redirected trolley in Loop exactly as he is in Switch.

\begin{thebibliography}{9}\small
\bibitem[Thomson(1985)]{Thomson1985} Thomson, J.~J. (1985). The trolley problem. \emph{The Yale Law Journal}, 94(6), 1395--1415.
\end{thebibliography}

\end{document}
